% ============================================================
%  IEEE Robotics and Automation Letters — Submission
%  "Image Verbalization for Semantic Room Analysis and
%   Autonomous Navigation on a 50-Gram Nano-Drone"
% ============================================================
\documentclass[journal]{IEEEtran}

\usepackage{cite}
\usepackage{amsmath,amssymb,amsfonts}
\usepackage{graphicx}
\usepackage{grffile}          % handles spaces and special chars in figure paths
\usepackage{textcomp}
\usepackage{xcolor}
\usepackage{booktabs}
\usepackage{multirow}
\usepackage{url}
\usepackage[hidelinks]{hyperref}
\usepackage[T1]{fontenc}

\newcommand{\todo}[1]{\textcolor{red}{\textbf{[TODO: #1]}}}

\begin{document}

\title{Image Verbalization for Semantic Room Analysis\\
and Autonomous Navigation on a 50-Gram Nano-Drone}

\author{Madhan~Kumar~B%
\thanks{M.~K.~B. is with the Department of Aerospace Engineering,
Indian Institute of Technology Madras, Chennai 600036, India
(e-mail: madanadan92@gmail.com).}}

\markboth{IEEE ROBOTICS AND AUTOMATION LETTERS}%
{Kumar: Image Verbalization for Nano-Drone Semantic Navigation}

\maketitle

% ============================================================
\begin{abstract}
Semantic room analysis and autonomous navigation on a 50-gram
nano-drone requires interpreting an environment from a
320$\times$240 camera frame --- only 85--102 image tokens per
LLM call.
We show that sub-VGA resolution is sufficient for safe autonomous
navigation, but only when each frame is enriched with structured
sensor metadata before the LLM call.
We call this process \emph{image verbalization}: YOLO-World
detections, DepthAnything\,v2 metric depth, and MediaPipe scene
flags are assembled with the raw JPEG into a structured prompt,
enabling a cloud multimodal LLM to simultaneously produce a
natural-language narration (streamed live to the operator) and a
navigation action executed autonomously via a modified R/C
transmitter.
A four-tier frequency architecture decouples the 2\,kHz PID
controller from the 0.1--0.4\,Hz LLM reasoner across four orders
of magnitude.
The key insight is that \textbf{action safety} --- not label
accuracy --- is the correct metric: a system that over-cautiously
hovers never advances into a hazard, even when its scene label is
wrong.
A 2,500$+$-trial parameter study establishes the jointly optimal
configuration (structured JSON, 256-token budget, 2-frame history,
$t=0$), overturning the ReAct $t=0.2$ convention for single-pass
classification.
A tier-ablation study proves each tier irreplaceable; a rule-based
baseline fails on sensor blindness; an LLM-without-metadata
baseline fails on depth-sensor errors.
Full-pipeline evaluation across 157 trials achieves
\textbf{97.5\%} action safety overall; Gemini\,2.5-Flash and
GPT-4o-mini reach \textbf{100\%}.
All four models achieve \textbf{100\% reason--action alignment}
in every trial.
Hardware validation on a live 50\,g custom drone across
\textbf{five real indoor rooms} confirms autonomous closed-loop
execution at \$0.17/hour.
\end{abstract}

\begin{IEEEkeywords}
Nano-drone autonomous navigation, multimodal large language models,
image verbalization, semantic scene analysis, indoor exploration,
low-resolution perception.
\end{IEEEkeywords}

% ============================================================
\section{Introduction}
\label{sec:intro}
% ============================================================

\IEEEPARstart{A}{utonomous} semantic navigation --- where a robot
enters an unknown room, understands what it sees, narrates the
scene aloud, and decides its next action --- is a long-standing
goal in robotics~\cite{ahn2022saycan,huang2022innermonologue}.
For nano-drones ($<$100\,g), the goal meets a hard constraint:
the onboard camera produces 320$\times$240 JPEG frames of only
9.7\,KB, consuming 85--102 image tokens per API call.

Does sub-VGA resolution prevent reliable semantic understanding?
\textbf{We show it does not} --- provided each frame is enriched
with structured sensor metadata before the LLM call.
This letter makes that enrichment process precise, calls it
\emph{image verbalization}, and validates it across 2,500$+$
trials and five real indoor rooms.

\textbf{The key insight is metric choice.}
Label accuracy (did the LLM classify the scene correctly?) is the
wrong optimisation target for autonomous navigation.
A model that upgrades a navigable scene to \texttt{caution} scores
poorly on label accuracy but produces a brief hover --- the correct
conservative behaviour.
A model that classifies a wall at 20\,cm as \texttt{safe} produces
a dangerous advance.
These errors are not equivalent.
\textbf{Action safety on hazard scenes} is the right target, and
every design decision is evaluated against it.

Three failure modes arise at 320$\times$240:
(1)~\textit{feature ambiguity} --- a grey wall and a grey doorway
are visually indistinguishable;
(2)~\textit{context loss} --- a person at 3\,m occupies
$\sim$20 pixels, defeating standard
detectors~\cite{persondetectionuav2025,kim2024yoloihd};
(3)~\textit{blocked-sensor blindness} --- zero detections from a
covered lens are misread as ``navigable space'' by any rule-based
engine.
Image verbalization addresses all three.

\textbf{Contributions:}
\begin{enumerate}
\item A systematic parameter study (2,500$+$ trials): jointly
  optimal LLM call configuration for semantic room analysis at
  320$\times$240. The ReAct $t=0.2$ convention is empirically
  overturned for single-pass classification.
\item A tier-ablation study proving each component individually
  irreplaceable, with hybrid MediaPipe triggering eliminating
  scene-change detection lag.
\item Full-pipeline evaluation (157 trials): 97.5\% action safety
  across four commercial LLMs (100\% for Gemini and GPT-4o-mini),
  all with 100\% reason--action alignment.
\item Closed-loop hardware demonstration on a live 50\,g custom
  drone across five real indoor rooms with a modified R/C
  transmitter executing autonomous roll/pitch commands from LLM
  decisions.
\end{enumerate}

% ============================================================
\section{Related Work}
\label{sec:related}
% ============================================================

\textbf{LLMs for robot planning.}
SayCan~\cite{ahn2022saycan} grounds LLM plans in structured robot
affordances; Inner Monologue~\cite{huang2022innermonologue} closes
the language feedback loop.
Both bypass raw visual input.
ChatGPT for Robotics~\cite{vemprala2023chatgpt} generates motion
primitives but does not address live semantic narration from a
flying platform.

\textbf{Vision-language models at low resolution.}
GPT-4V~\cite{achiam2023gpt4}, Gemini~\cite{reid2024gemini}, and
LLaVA~\cite{liu2023llava} are benchmarked at $\geq$720p.
A key finding of this letter is that CLIP-ViT features, effective
at higher resolutions, add noise at 320$\times$240 and should be
removed.

\textbf{UAV perception.}
Sub-VGA UAV-perspective person detection degrades
severely~\cite{persondetectionuav2025,kim2024yoloihd}.
Recent search-and-rescue~\cite{rescuedrone2025} and hierarchical
control~\cite{hierarchicaldrone2025} systems operate on
higher-resolution feeds without the strict 10\,KB frame-size
constraint of a weight-limited nano-drone.

\textbf{Monocular depth.}
DepthAnything\,v2~\cite{yang2024depthanythingv2} systematically
overestimates distance on featureless surfaces at close range.
Section~\ref{sec:full_pipeline} quantifies the navigation
consequence and shows how explicit reliability metadata allows the
LLM to override erroneous readings.

\textbf{Temperature convention.}
ReAct~\cite{yao2022react} establishes $t=0.2$ for iterative
agentic tasks.
We show that for single-pass scene classification, $t=0$ is
strictly superior (Section~\ref{sec:temp}).

% ============================================================
\section{Image Verbalization Architecture}
\label{sec:arch}
% ============================================================

\subsection{Core Concept}

A 320$\times$240 frame alone is often semantically ambiguous.
Paired with structured sensor metadata, the same frame provides
sufficient grounding for the LLM to narrate the scene and decide
a navigation action.
This enrichment process --- \emph{image verbalization} --- proceeds
in three steps:
\begin{enumerate}
  \item Local detectors (Tiers\,1.5 and 2) extract object classes,
    bounding boxes, metric depth, and scene flags in $\leq$275\,ms,
    with zero API cost;
  \item A structured prompt assembles: JPEG\,+\,metadata\,+\,
    2-frame context history;
  \item The LLM returns structured JSON:
    \texttt{scene\_description} (narrated live to operator) and
    \texttt{recommended\_action} (executed autonomously).
\end{enumerate}

\subsection{Four-Tier Frequency Architecture}

Each function is assigned to the fastest tier its computational
budget permits (Fig.~\ref{fig:concept}):

\begin{description}
\item[Tier\,1 --- PID Firmware (2\,kHz):]
  Attitude, altitude, and position hold on the ESP32-S3.
  Completely independent of the LLM; never stalled by an API call.

\item[Tier\,1.5 --- MediaPipe Scene Monitor ($\sim$60\,Hz):]
  EfficientDet-Lite0~\cite{tan2020efficientdet,lugaresi2019mediapipe}
  performs three binary checks in $\sim$16\,ms at zero API cost:
  person confidence\,$>$\,0.30, wall-fill gradient standard
  deviation\,$<$\,30, mean luminance\,$<$\,40.
  Any positive flag fires an immediate Tier\,3 LLM call, bypassing
  the scheduled polling interval.

\item[Tier\,2 --- Structural Perception ($\sim$4\,Hz):]
  YOLO-World~\cite{cheng2024yoloworld} (17 indoor hazard classes)
  and DepthAnything\,v2 Metric
  Indoor~\cite{yang2024depthanythingv2} run in $\sim$275\,ms and
  produce the sensor metadata string for every LLM prompt.
  CLAHE pre-processing~\cite{clahe_zuiderveld1994}
  (\texttt{clip\_limit}=2.0, \texttt{tile\_grid}=$(8,8)$) is
  applied before all inference.

\item[Tier\,3 --- LLM Cognitive Reasoner (0.1--0.4\,Hz):]
  Gemini\,2.5-Flash in production~\cite{reid2024gemini}.
  Structured JSON output enforces separate \texttt{risk\_level}
  and \texttt{recommended\_action} fields, making contradictory
  outputs detectable as parse errors.
\end{description}

\begin{figure}[t]
  \centering
  \includegraphics[width=\columnwidth]{Thesis Figures/Chapter 2 Architecture.png}
  \caption{Four-tier image verbalization architecture.
    A 320$\times$240 JPEG (85--102 image tokens) is enriched by
    Tier\,1.5 and Tier\,2 before the Tier\,3 LLM call.
    The LLM produces a live verbal narration \emph{and} a
    navigation action in a single call.}
  \label{fig:concept}
\end{figure}

% ============================================================
\section{Experimental Setup}
\label{sec:setup}
% ============================================================

All experiments use eight canonical scenes
(Table~\ref{tab:scenes}), captured from the ESP32-S3-Sense at
320$\times$240 with five independent repetitions
(40 frames/scene, 320 total).
Four commercial multimodal LLMs are evaluated using the HELM
methodology~\cite{liang2023helm}, five independent runs per
condition.

Primary metrics: \textbf{action safety} (ActSafe --- recommended
action does not advance into a ground-truth hazard) and
\textbf{dangerous cases} (truth=hazard and
action=\textsc{pitch\_forward}).
Label accuracy (LblAcc) is secondary; it penalises conservative
over-caution equally with genuine failure.

\begin{table}[t]
  \caption{Eight canonical scenes. Ground-truth:
    \texttt{hazard}$\to$stop/hover;
    \texttt{caution}$\to$hover; \texttt{safe}$\to$advance.}
  \label{tab:scenes}
  \centering
  \small
  \begin{tabular}{lll}
    \toprule
    Scene & Label & Principal challenge \\
    \midrule
    \texttt{person\_near}  & hazard  & Person $<$0.5\,m; priority escalation \\
    \texttt{wall\_close}   & hazard  & DA\,v2 error: 2.09\,m reported at 20\,cm \\
    \texttt{blocked\_lens} & hazard  & Zero detections from occluded lens \\
    \texttt{person\_far}   & caution & Person $>$2\,m; over-caution expected \\
    \texttt{dim\_light}    & caution & Low luminance ($\approx$40--80) \\
    \texttt{cluttered}     & caution & Many objects, no dominant block \\
    \texttt{object\_table} & caution & Lower-frame partial obstruction \\
    \texttt{door\_open}    & safe    & Open passage; must advance \\
    \bottomrule
  \end{tabular}
\end{table}

% ============================================================
\section{Validation: LLM Call Parameter Optimisation}
\label{sec:validation}
% ============================================================

Single-variable experiments isolate each LLM call parameter
(8 scenes $\times$ 5 runs per condition).
Table~\ref{tab:validation} summarises all results.

\subsection{Baseline Model Comparison}

160 trials across four models with a combined
YOLO\,+\,MediaPipe\,+\,CLIP metadata prompt.
All four models produce zero dangerous cases from 320$\times$240
frames, confirming core feasibility: multimodal LLMs can reason
safely at 85--102 image tokens per call.
Gemini leads on latency (2,468\,ms); GPT-4o leads on label
accuracy (70\%) and narration quality (4.4/5).

\subsection{Output Token Budget and CLIP Ablation}

Narration quality saturates at 256 tokens ($\Delta Q\!=\!0.01$
from 256 to 512).
CLIP-ViT changes quality by $\leq$0.05 at all budget levels ---
its scores fall within $\pm$0.013 of the uniform 0.200 baseline
across all scenes: no scene-discriminative signal at 320$\times$240.
In one configuration (Gemini at 128 tokens), removing CLIP
\emph{improves} quality by $+$0.40/5 because the uninformative
label consumes constrained reasoning capacity.

\textbf{Decision: \texttt{max\_tokens=256}; CLIP removed.}

\subsection{Sampling Temperature}
\label{sec:temp}

Temperature swept over $[0.0,\,0.2,\,0.5,\,0.8,\,1.0]$ across
800 trials.
$t=0$ achieves highest label accuracy (61.3\%), zero dangerous
cases, and the lowest flip rate (14.1\%).
Dangerous cases appear only at $t\!\geq\!0.5$, exclusively
GPT-4o-mini on \texttt{wall\_close}.
Reasoning quality is temperature-invariant ($\approx$94\%
reasoning-risk alignment regardless of $t$): degradation above
$t=0.5$ is a consistency failure, not a reasoning failure.

\textbf{Decision: $t=0$, overturning ReAct
$t=0.2$~\cite{yao2022react} for single-pass classification.}

\subsection{Prompt Technique}

Five techniques across $\approx$800 trials.
Structured JSON achieves the highest label accuracy (56.2\%) and
zero dangerous cases.
The schema enforces separate \texttt{risk\_level} and
\texttt{recommended\_action} fields --- contradictory outputs are
detectable as parse errors, architecturally preventing the most
dangerous failure mode.
Few-shot-3 produces 5 dangerous cases (all Gemini,
\texttt{wall\_close}).
CoT-300 is truncated 67\% of the time, collapsing action reasoning
presence to 32.5\%.

\textbf{Decision: Structured JSON adopted.}

\subsection{Scene Context History}

Three history modes across 300 trials.
Short history (2 frames, $+$38 tokens) improves label accuracy by
16.2\,pp over stateless ($43.0\%\!\to\!59.2\%$).
GPT-4o-mini stateless achieves only 4\% label accuracy (near
chance), rising to 48\% with 2-frame history.
All modes maintain 100\% action safety and zero dangerous cases.
The system exhibits a desirable asymmetry: hazard onset detected
100\%; hazard clearance only 26.7\% --- the drone lingers in
\textsc{hover} rather than prematurely resuming forward motion.

\textbf{Decision: 2-frame short history.
Final validated configuration: structured JSON,
\texttt{max\_tokens=256}, 2-frame history, $t=0$, no CLIP.}

\begin{table}[t]
  \caption{Validation experiment results.
    (a)~Token budget. (b)~Temperature (800 trials).
    (c)~Prompt technique ($\approx$800 trials).
    (d)~Context history (300 trials).}
  \label{tab:validation}
  \centering
  \small
  \textbf{(a) Token budget}\vspace{2pt}\\
  \begin{tabular}{rrrr}
    \toprule
    \texttt{max\_tokens} & Quality/5 & LblAcc & Latency \\
    \midrule
    64           & 3.27 & 36.7\% & 3,142\,ms \\
    128          & 4.15 & \textbf{58.8\%} & 3,649\,ms \\
    \textbf{256} & \textbf{4.48} & 57.3\% & 4,228\,ms \\
    512          & 4.49 & 57.7\% & 4,294\,ms \\
    \bottomrule
  \end{tabular}\vspace{4pt}

  \textbf{(b) Temperature}\vspace{2pt}\\
  \begin{tabular}{crrr}
    \toprule
    Temp & LblAcc & ActSafe & FlipRate \\
    \midrule
    \textbf{0.0} & \textbf{61.3\%} & \textbf{100\%} & \textbf{14.1\%} \\
    0.2 & 59.4\% & 100\%  & 21.9\% \\
    0.5 & 51.6\% & 98.7\% & 30.7\% \\
    0.8 & 58.1\% & 98.8\% & 20.3\% \\
    1.0 & 56.6\% & 99.4\% & 32.8\% \\
    \bottomrule
  \end{tabular}\vspace{4pt}

  \textbf{(c) Prompt technique}\vspace{2pt}\\
  \begin{tabular}{lrrr}
    \toprule
    Technique & LblAcc & ActSafe & Dangerous \\
    \midrule
    Zero-shot              & 41.9\% & \textbf{100\%} & \textbf{0} \\
    Few-shot-3             & 51.2\% & 96.9\%         & 5 \\
    CoT-300 (67\% trunc.)  & 18.8\% & \textbf{100\%} & \textbf{0} \\
    CoT-600                & 46.8\% & 98.8\%         & 0 \\
    \textbf{Struct.\ JSON} & \textbf{56.2\%} & \textbf{100\%} & \textbf{0} \\
    \bottomrule
  \end{tabular}\vspace{4pt}

  \textbf{(d) Context history}\vspace{2pt}\\
  \begin{tabular}{lrrrr}
    \toprule
    Mode & LblAcc & ChgDetect & ActSafe & $+$tokens \\
    \midrule
    Stateless              & 43.0\% & 65.0\% & \textbf{100\%} & --- \\
    \textbf{Short (2 fr.)} & \textbf{59.2\%} & \textbf{67.5\%} & \textbf{100\%} & $+$38 \\
    Full (all fr.)         & 50.5\% & 56.4\% & \textbf{100\%} & $+$59 \\
    \bottomrule
  \end{tabular}
\end{table}

% ============================================================
\section{Integration: Proof of the Layered Architecture}
\label{sec:integration}
% ============================================================

\subsection{Timescale Separation}

Table~\ref{tab:timescale} and Fig.~\ref{fig:latency} show all tier
latencies measured across five runs and 160 LLM calls.
The four tiers span four orders of magnitude.
Running Tier\,3 at Tier\,2's frequency (3--4\,Hz) would increase
Gemini API cost by $\sim$10$\times$/minute at no safety benefit.
Non-overlapping 95\% CIs between Claude and Gemini confirm the
3.1$\times$ latency difference is statistically significant.

\begin{table}[t]
  \caption{Measured tier latencies (mean, 5 runs, 160 LLM calls).}
  \label{tab:timescale}
  \centering
  \small
  \begin{tabular}{llrl}
    \toprule
    Tier & Component & Mean latency & Frequency \\
    \midrule
    1   & PID controller         & 0.5\,ms   & 2,000\,Hz \\
    1.5 & MediaPipe EfficientDet & 16.2\,ms  & $\sim$60\,Hz \\
    2   & YOLO-World + DA\,v2    & 274.5\,ms & $\sim$4\,Hz \\
    3   & Gemini\,2.5-Flash      & 2,333\,ms & $\sim$0.4\,Hz \\
    3   & GPT-4o                 & 2,627\,ms & $\sim$0.4\,Hz \\
    3   & GPT-4o-mini            & 4,333\,ms & $\sim$0.2\,Hz \\
    3   & Claude (Sonnet)        & 7,218\,ms & $\sim$0.1\,Hz \\
    \bottomrule
  \end{tabular}
\end{table}

\begin{figure}[t]
  \centering
  \includegraphics[width=\columnwidth]{Thesis Figures/G4 — fig_4 Latency Analysis.pdf}
  \caption{Four-tier latency cascade and LLM latency comparison
    with 95\% CI. Adjacent tiers differ by at least one order of
    magnitude; collapsing any two would impose API latency on a
    safety-critical layer or prohibitive cost on a
    high-frequency layer.}
  \label{fig:latency}
\end{figure}

\subsection{Tier Isolation and Baseline Comparison}
\label{sec:tier_isolation}

Four conditions are compared (Table~\ref{tab:isolation}).

\textbf{Rule-based baseline (Tier\,2-only):}
YOLO-World sees zero objects from a covered or dark lens and
returns \textsc{pitch\_forward}: 5/5 dangerous cases on
\texttt{blocked\_lens}.
No rule engine can detect that its own sensor has been compromised.

\textbf{LLM-without-metadata baseline:}
GPT-4o-mini without Tier\,2 depth metadata produces 5/5 dangerous
cases on \texttt{wall\_close}: without metric depth, the
featureless grey wall at 20\,cm is indistinguishable from clear
space at 320$\times$240.
With depth metadata added, this drops to 1/5.

\textbf{Full pipeline:}
Gemini (full pipeline) achieves zero dangerous cases across all
conditions.
All LLMs recognise occluded or dark frames directly from image
pixels --- recovering the blocked-lens failure of the rule-based
baseline.

\begin{table}[t]
  \caption{Tier isolation and baseline comparison
    (40 trials per condition, 8 scenes $\times$ 5 runs).}
  \label{tab:isolation}
  \centering
  \small
  \begin{tabular}{p{3.0cm}lrr}
    \toprule
    Configuration & Model & ActSafe & Dangerous \\
    \midrule
    \textit{Baseline:} Rule-based   & ---         & 87.5\% & 5 \\
    \textit{Baseline:} LLM, no meta & GPT-4o-mini & 87.5\% & 5 \\
    \midrule
    Tier\,1.5 alone  & Rule-based  & \textbf{100\%} & 0 \\
    Tier\,3 full     & Gemini      & \textbf{100\%} & 0 \\
    Tier\,3 full     & GPT-4o      & \textbf{100\%} & 0 \\
    Full pipeline    & Gemini      & \textbf{100\%} & 0 \\
    Full pipeline    & GPT-4o      & 95.0\%         & 2 \\
    Full pipeline    & Claude      & 97.4\%         & 1 \\
    Full pipeline    & GPT-4o-mini & 97.5\%         & 1 \\
    \bottomrule
  \end{tabular}
\end{table}

\subsection{Hybrid Trigger}

Scheduled polling detects a hazard 3 frames late (one polling
interval at 30\,fps).
Hybrid triggering fires an additional LLM call immediately when
MediaPipe flags a positive check --- detecting at the exact entry
frame (0 frames late).
At 0.5\,m/s, 3 frames is $\sim$5\,cm of uninformed forward travel.
In one run, YOLO-World misclassified a kneeling person as a
structural wall (confidence 0.25); MediaPipe triggered the LLM
correctly (person confidence 0.391), and Gemini's narration:
\emph{``my visual analysis confirms they are much closer than the
YOLO detection suggested''} --- demonstrating active
sensor-disagreement resolution in live verbal output.
Both strategies achieve 100\% hazard catch rate; the difference is
timing only.
Cost: $+$\$0.00012 per hazard event (Gemini).

\textbf{Hybrid triggering adopted.}

\subsection{Full Pipeline Evaluation}
\label{sec:full_pipeline}

The complete pipeline (structured JSON, \texttt{max\_tokens=256},
no CLIP, 2-frame history, $t=0$, hybrid trigger) across 157 valid
trials (Table~\ref{tab:g5}, Fig.~\ref{fig:g5}).

\textbf{Gemini and GPT-4o-mini achieve 100\% action safety.}

GPT-4o's 3 dangerous cases all arise on \texttt{wall\_close}:
DepthAnything\,v2 reports 2.09\,m for a wall at 20\,cm, and
GPT-4o anchors on the number rather than the visual texture fill.
A sensor-reliability metadata flag reduced dangerous cases by 79\%
between run~1 and run~2 (14\,$\to$\,3).
Gemini was 100\% error-free before and after the flag:
\emph{``The image shows a uniformly light gray surface\ldots
The depth\_m of 2.09\,m seems inconsistent with the visual
evidence\ldots Risk: hazard. Action: HOVER.''} --- resolving the
sensor error from image pixels alone.

\textbf{All four models achieve 100\% reason--action alignment:}
no model states a hazard classification and recommends
\textsc{pitch\_forward} in any trial.
GPT-4o's 3 errors are internally consistent --- wrong depth
perception $\to$ correct reasoning from wrong input $\to$
consistent but dangerous action.
The failure is input perception, not LLM reasoning.

The per-scene breakdown (Table~\ref{tab:g5}b) reveals the
action-safety insight: \texttt{person\_far} achieves only 11\%
label accuracy but 100\% action safety --- all errors are
conservative \textsc{hover}s, never dangerous advances.
Only \texttt{wall\_close}/GPT-4o produces genuinely unsafe actions.

\begin{table}[t]
  \caption{Full pipeline evaluation (157 valid trials).
    (a)~Per-model. (b)~Per-scene (all models combined).}
  \label{tab:g5}
  \centering
  \small
  \textbf{(a) Per-model}\vspace{2pt}\\
  \begin{tabular}{lrrrr}
    \toprule
    Model & $N$ & LblAcc & ActSafe & Dangerous \\
    \midrule
    Claude              & 37  & 54.0\% & 97.0\%          & 0 \\
    GPT-4o              & 40  & 52.0\% & 92.5\%          & 3 \\
    GPT-4o-mini         & 40  & \textbf{80.0\%} & \textbf{100\%} & 0 \\
    \textbf{Gemini}     & 40  & 50.0\% & \textbf{100\%}  & 0 \\
    \midrule
    \textbf{Overall}    & 157 & 59.2\% & \textbf{97.5\%} & 3 \\
    \bottomrule
  \end{tabular}
  \vspace{4pt}

  \textbf{(b) Per-scene}\vspace{2pt}\\
  \begin{tabular}{llrr}
    \toprule
    Scene & Label & LblAcc & ActSafe \\
    \midrule
    \texttt{person\_near}  & hazard  & 95\% & \textbf{100\%} \\
    \texttt{wall\_close}   & hazard  & 80\% & 85\%$^\dagger$ \\
    \texttt{blocked\_lens} & hazard  & 75\% & 95\% \\
    \texttt{door\_open}    & safe    & 75\% & \textbf{100\%} \\
    \texttt{dim\_light}    & caution & 50\% & \textbf{100\%} \\
    \texttt{object\_table} & caution & 50\% & \textbf{100\%} \\
    \texttt{cluttered}     & caution & 35\% & \textbf{100\%} \\
    \texttt{person\_far}   & caution & 11\% & \textbf{100\%} \\
    \bottomrule
    \multicolumn{4}{l}{$^\dagger$ GPT-4o only; DA\,v2 depth anchor.}
  \end{tabular}
\end{table}

\begin{figure}[t]
  \centering
  \includegraphics[width=\columnwidth]{Thesis Figures/G5 — fig_7 Real Vision Pipeline Analysis.pdf}
  \caption{Full pipeline analysis.
    (a)~Label accuracy vs.\ action safety per model: GPT-4o-mini
    has 80\% label accuracy and 100\% action safety; its errors
    are conservative \textsc{hover}s.
    (b)~Action safety heatmap; only \texttt{wall\_close}/GPT-4o
    cell is unsafe.
    (c)~Gemini overrides DA\,v2 depth error from image pixels.
    (d)~Before/after sensor-reliability flag: 79\% reduction
    in dangerous cases.}
  \label{fig:g5}
\end{figure}

% ============================================================
\section{Hardware Validation: Five Real Rooms}
\label{sec:h5}
% ============================================================

\subsection{Platform}

Custom 50\,g airframe with bespoke flight controller on the
ESP32-S3-Sense.
Tier\,1 PID runs at 2\,kHz on-chip, fully independent of all
cloud API calls.
Tiers\,1.5--3 run asynchronously on an off-board companion
computer over Wi-Fi; no API call ever stalls the flight control
loop.

\subsection{Transmitter Modification}

To execute LLM navigation decisions without modifying flight
controller firmware, a standard R/C transmitter's roll and pitch
signal lines are intercepted by a microcontroller that injects
PWM offsets corresponding to the LLM action:
\begin{itemize}
  \item \textsc{pitch\_forward}: $+\delta_\text{pitch}$
    (calibrated for 0.5\,m/s indoor cruise);
  \item \textsc{hover}: zero offset on both axes;
  \item \textsc{lateral\_left/right}: $\pm\delta_\text{roll}$.
\end{itemize}
The Tier\,1 PID receives only the resulting setpoints and remains
unaware of the LLM.
The operator retains physical override at all times.

\subsection{Five-Room Protocol and Results}

The complete pipeline is deployed in five real indoor rooms:
Office (clutter, natural daylight), Corridor (narrow passage,
person encounter), Storage (dim lighting, occluded lens event),
Lab (open area, dynamic obstacles), and Foyer (door/passage
decision, mixed lighting).
Four hazard events are injected per room (3 repetitions each):
\texttt{person\_near}, \texttt{wall\_close},
\texttt{blocked\_lens}, \texttt{dim\_light}.
60 triggered LLM calls per room, 300 total.

\begin{table}[t]
  \caption{Hardware validation across five real indoor rooms
    (300 triggered LLM calls total, Gemini\,2.5-Flash).
    \todo{Fill with measured data.}}
  \label{tab:h5}
  \centering
  \small
  \begin{tabular}{lrrrr}
    \toprule
    Room & E2E latency & ActSafe & Quality/5 & Cost/min \\
    \midrule
    Office (R1)   & \todo{} & \todo{} & \todo{} & \todo{} \\
    Corridor (R2) & \todo{} & \todo{} & \todo{} & \todo{} \\
    Storage (R3)  & \todo{} & \todo{} & \todo{} & \todo{} \\
    Lab (R4)      & \todo{} & \todo{} & \todo{} & \todo{} \\
    Foyer (R5)    & \todo{} & \todo{} & \todo{} & \todo{} \\
    \midrule
    Overall       & \todo{} & \todo{} & \todo{} & \todo{} \\
    \bottomrule
  \end{tabular}
\end{table}

% ============================================================
\section{Discussion}
\label{sec:discussion}
% ============================================================

\textbf{Low resolution is sufficient with metadata enrichment.}
YOLO-World labels and DepthAnything\,v2 metric depth convert a
visually ambiguous sub-VGA frame into a semantically grounded
prompt.
CLIP-ViT provides no discriminative signal at 320$\times$240 and
degrades quality under tight token budgets --- a counterintuitive
finding for practitioners who assume more modalities always help.

\textbf{Action safety is the right metric.}
\texttt{person\_far} at 11\% label accuracy but 100\% action
safety is the clearest demonstration: all labelling errors produce
\textsc{hover}, never a dangerous advance.
Designing for action safety on hazard scenes is the correct
engineering objective for autonomous exploration drones.

\textbf{LLMs actively resolve sensor errors.}
Gemini's verbatim response on \texttt{wall\_close} demonstrates
reasoning about sensor-metadata inconsistency from image pixels
alone, without the sensor-reliability flag.
The flag reduced GPT-4o dangerous cases by 79\%, establishing a
design principle: give the LLM explicit reliability context for
noisy sensor suites.

\textbf{Limitations.}
LLM latency (2--7\,s) limits navigable speed to $\leq$0.5\,m/s;
faster platforms require edge-deployed LLMs.
The eight offline scenes are laboratory-controlled; H5
five-room deployment tests generalisability but does not cover
outdoor or highly dynamic environments.

% ============================================================
\section{Conclusion}
\label{sec:conclusion}
% ============================================================

We presented image verbalization: a four-tier architecture that
enables a 50-gram nano-drone to enter an unknown indoor room,
produce a live natural-language narration, and execute autonomous
navigation decisions from a 9.7\,KB, 320$\times$240 camera frame
at \$0.17/hour.
Despite only 85--102 image tokens per call, the full pipeline
achieves 97.5\% action safety across four commercial LLMs, with
Gemini and GPT-4o-mini at 100\%.
A tier-ablation study proved each component irreplaceable: the
rule-based baseline fails on sensor blindness; the
LLM-without-metadata baseline fails on depth-sensor errors; only
the full four-tier pipeline handles both.
A systematic parameter study across 2,500$+$ trials established
the optimal configuration and overturned the ReAct $t=0.2$
convention for single-pass classification.
Hardware validation across five real indoor rooms confirmed
autonomous closed-loop execution with full decoupling of the
2\,kHz PID from cloud API latency.

% ============================================================
\appendix
\section{Production Configuration}
\label{sec:config}

\begin{table}[t]
  \caption{Production configuration. Each parameter established
    by a specific validation or integration experiment.}
  \label{tab:config}
  \centering
  \small
  \begin{tabular}{lll}
    \toprule
    Parameter & Value & Established by \\
    \midrule
    Prompt       & Structured JSON              & Prompt technique \\
    Token budget & \texttt{max\_tokens=256}     & Token budget \\
    CLIP         & Removed                      & Token budget \\
    History      & 2-frame short ($+$38 tokens) & Context history \\
    Temperature  & $t=0.0$                      & Temperature \\
    Architecture & Four-tier                    & Timescale separation \\
    Trigger      & Hybrid (sched.+MediaPipe)    & Hybrid trigger \\
    LLM          & Gemini\,2.5-Flash            & Full pipeline \\
    \bottomrule
  \end{tabular}
\end{table}

% ============================================================
\bibliographystyle{IEEEtran}
\bibliography{Thesis/references}

\end{document}
